Short-term memory can guide navigation by limiting repeated visits to recently explored locations. Neuronal traces describe this persistence mathematically, but do not uniquely identify its cellular substrate. Astrocytes are a candidate because their calcium-dependent interactions with neurons can influence neuronal activity and memory. Their calcium dynamics could therefore provide a transient record of experience that biases subsequent choices. We examine this idea in AstroRL, a recurrent spiking network with astrocyte-inspired inhibitory regulation. Reduced glutamate-calcium dynamics retain recent spatial inputs, while learned predictions recruit inhibition for actions with familiar predicted outcomes. Prediction error trains the recurrent network, and reward trains a detached motor head. In the tested navigation tasks, memory increases exploration without policy learning and improves goal-directed performance in structured environments after reinforcement learning. Disabling memory in trained controllers reduces performance. The astrocyte-inspired memory module performed competitively with neuronal traces, providing computational support for the hypothesis that astrocytic calcium dynamics could fulfill analogous memory functions in biological circuits. The results illustrate how persistent astrocyte-inspired regulation of inhibition can connect recent spatial experience to action selection.